\section{Network Time (NTP)}

\subsection{Overview}

Time synchronization is one of the most critical and least-visible aspects of network management. All lab services -- Kerberos authentication, TLS certificate validation, log correlation, distributed databases, and filesystem consistency -- depend on time accuracy. The lab uses a two-tier NTP architecture:

\begin{enumerate}
    \item \textbf{Stratum 1 source}: GPS-referenced clock on the \texttt{time} host (10.10.12.2) provides the lab's authoritative time reference.
    
    \item \textbf{Stratum 2+ clients}: All other hosts sync to \texttt{time.lab.bitsmasher.net}.
\end{enumerate}

\subsection{The Time Host}

\begin{itemize}
    \item \textbf{Hostname}: time.lab.bitsmasher.net (alias: time)
    
    \item \textbf{IP}: 10.10.12.2
    
    \item \textbf{User}: openclaw
    
    \item \textbf{OS}: Debian 12 bookworm
    
    \item \textbf{Role}: Stratum-1 NTP server / reference clock source for the lab
\end{itemize}

The time host has a GPS unit connected (serial/tty interface) providing UTC discipline. Its ntpd is configured via ntpsec to serve the lab subnet and use the GPS PPS signal as the primary time reference.

\subsection{Configuration on Client Hosts}

The NTP configuration is managed by Ansible via the \texttt{ntp} role in the lab-franklin collection. Here is what a typical client configuration looks like (as seen on stargate):

\begin{lstlisting}[style=mystyle]
# /etc/ntpsec/ntp.conf -- managed by Ansible
driftfile /var/lib/ntp/ntp.drift
leapfile /usr/share/zoneinfo/leap-seconds.list

statsdir /var/log/ntpstats/

statistics loopstats peerstats clockstats
filegen loopstats file loopstats type day enable
filegen peerstats file peerstats type day enable
filegen clockstats file clockstats type day enable

server time.lab.bitsmasher.net
restrict 10.10.8.0/21
restrict -4 default
\end{lstlisting}

Key points:
\begin{itemize}
    \item \texttt{driftfile}: Tracks oscillator drift between NTP restarts. Critical for stability on hosts that reboot frequently (Jetson devices, test VMs).
    
    \item \texttt{statistics}: Logs loop/peer/clock stats daily in /var/log/ntpstats/. Use these to verify sync health.
    
    \item \texttt{server}: Points to the lab's time host -- never pool.ntp.org for infrastructure hosts that need deterministic time sources.
    
    \item \texttt{restrict}: 
    \begin{itemize}
        \item \texttt{10.10.8.0/21} grants query access to the main lab subnet (used by other internal clients or monitoring tools).
        
        \item \texttt{-4 default} blocks all IPv4 queries not explicitly allowed, providing a deny-by-default posture.
    \end{itemize}
\end{itemize}

\subsection{Installation and Setup Steps}

For a new host that needs NTP:

\begin{enumerate}
    \item Install ntpsec packages:
    \begin{lstlisting}[style=mystyle]
apt update && apt install -y ntpsec
    \end{lstlisting}

    \item Configure /etc/ntpsec/ntp.conf (either manually or via Ansible ntp role):
    \begin{enumerate}
        \item Set \texttt{server time.lab.bitsmasher.net} as the upstream reference.
        
        \item Set timezone: \texttt{timedatectl set-timezone America/Denver} (or appropriate TZ).
        
        \item Create the log directory if it doesn't exist: \texttt{mkdir -p /var/log/ntpstats}.
    \end{enumerate}

    \item Start and enable the service:
    \begin{lstlisting}[style=mystyle]
systemctl enable --now ntpsec.service
timedatectl  # verify "System clock synchronized: yes"
    \end{lstlisting}
\end{enumerate}

\subsection{Verification on Clients}

Check synchronization status with these commands:

\begin{lstlisting}[style=mystyle]
# Check if the system clock is synchronized
timedatectl | grep "System clock synchronized"

# Check NTP peer status (if ntpq is available)
ntpq -p

# View recent sync logs
journalctl -u ntpsec --since today
    \end{lstlisting}

\subsection{Notes for Deployment}

\begin{itemize}
    \item The Ansible \texttt{ntp} role handles all of the above automatically across lab hosts. It installs packages, copies configuration, sets timezone, and creates log directories.
    
    \item For Jetson devices (node900--node903), drift compensation is especially important -- they boot frequently and have less stable oscillators than desktop/server hardware.
    
    \item The GPS unit on the time host should be verified periodically: confirm it's locking to satellites, not using its internal oscillator as a holdover source. Check this by running \texttt{ntptime} on the time host itself.
    
    \item If the time host goes offline and all clients lose sync beyond their maximum correction threshold (typically 128 seconds), systems may need manual time correction before Kerberos/TLS will function again.
\end{itemize}
